\section{作业答案速查}
\subsection*{第1周}
\begin{enumerate}
\item \hyperref[ex:w01-t1]{教材第9--10页第1题}：依次为$A\cup B\cup C$，$AB\cup AC\cup BC$，$A^cB^cC^c$，$(AB\cup AC\cup BC)^c$，$(ABC)^c$，$ABC^c\cup AB^cC\cup A^cBC$。
\item \hyperref[ex:w01-t2]{教材第10页第2题}：$13/30$。
\item \hyperref[ex:w01-t4]{教材第10页第4题}：有放回$169/512$；无放回$3/8$。
\item \hyperref[ex:w01-t5]{教材第10页第5题}：$p_1=[m(m-1)+n(n-1)]/[M(M-1)]$，$p_2=(m^2+n^2)/M^2$，且$p_2-p_1=2mn/[M^2(M-1)]>0$。
\item \hyperref[ex:w01-t6]{教材第10页第6题}：$1/2$。
\item \hyperref[ex:w01-t7]{教材第10页第7题}：按原题“小于”理解，$0<a<10/9$，不存在正的最小边长。
\item \hyperref[ex:w01-h2]{作业第2题：可列德摩根律}：两式均成立。
\item \hyperref[ex:w01-h3]{作业第3题：恰有两个事件发生}：两边均表示恰有两个事件发生。
\end{enumerate}
\subsection*{第2周}
\begin{enumerate}
\item \hyperref[ex:w02-t1]{教材第19页第1题}：不存在可数无穷等可能且每个结果有正概率的模型。
\item \hyperref[ex:w02-t2]{教材第19页第2题(1)}：$\mathscr E$是$E$上的$\sigma$代数，补运算必须相对于$E$进行。
\item \hyperref[ex:w02-t4]{教材第19页第4题}：所给下界成立；结合非负性还可写成$P(AB)\geq\max\{0,P(A)+P(B)-1\}$。
\item \hyperref[ex:w02-t5]{教材第19页第5题}：$P(A\cap B)\in[1/12,1/3]$。
\item \hyperref[ex:w02-h2]{作业第2题：偏离独立乘积的界}：$1/4$为可达到的界，例如$A=B$且$P(A)=1/2$时取等号。
\item \hyperref[ex:w02-h3]{作业第3题：开区间的可列表示}：该可列并表示成立。
\item \hyperref[ex:w02-h4]{作业第4题：无限重复试验的事件极限}：单减；极限为永远不成功；至少成功一次的概率为1。
\end{enumerate}
\subsection*{第3周}
\begin{enumerate}
\item \hyperref[ex:w03-t3]{教材第19页第3题}：两项均成立，且$\mu(\Omega)=\sum_{k=1}^n\mu_k(\Omega_k)$。
\item \hyperref[ex:w03-t6]{教材第19页第6题}：可列个满概率事件的交仍为满概率事件。
\item \hyperref[ex:w03-t241]{教材第24页第1题}：结论成立，且题设自动保证$P(A)>0$。
\item \hyperref[ex:w03-t242]{教材第24页第2题}：(1)一定成立；(2)不一定，以上反例中$P(A)=3/10<P(B)=7/10$。
\item \hyperref[ex:w03-t254]{教材第25页第4题}：$528/5915\approx0.08926$。
\item \hyperref[ex:w03-h2]{作业第2题：一排夫妻均不相邻}：$\displaystyle\sum_{r=0}^n(-1)^r\binom nr2^r(2n-r)!/(2n)!$。
\item \hyperref[ex:w03-h3]{作业第3题：已知补事件条件概率}：$0.73$。
\item \hyperref[ex:w03-h4]{作业第4题：玻璃器皿连续三次未摔坏}：$0.028$。
\item \hyperref[ex:w03-h5]{作业第5题：两种天气的递推}：$\bigl[1+(2p-1)^{n-1}\bigr]/2$。
\end{enumerate}
\subsection*{第4周}
\begin{enumerate}
\item \hyperref[ex:w04-t25-3]{教材第25页第3题}：$\P(D)=(1+4p)/5$；$\P(K\mid D)=5p/(1+4p)$。
\item \hyperref[ex:w04-t25-5]{教材第25页第5题}：两种信号的条件误码率均为$0.05$时，答案为$38/39$；只给总误码率则条件不足。
\item \hyperref[ex:w04-t25-6]{教材第25页第6题}：$\displaystyle\frac{\sum_{i=1}^ki^{n+1}}{k\sum_{i=1}^ki^n}$。
\item \hyperref[ex:w04-t32-1]{教材第32页第1题}：两个结论均成立；事件与自身独立当且仅当其概率为0或1。
\item \hyperref[ex:w04-t32-2]{教材第32页第2题}：3孩时独立；4孩时不独立。
\item \hyperref[ex:w04-t33-6]{教材第33页第6题}：$p_n=\frac12+(c-\frac12)(2p-1)^{n-1}$，极限为$1/2$。
\item \hyperref[ex:w04-h3]{作业第3题：两个并事件的独立性}：$A_1\cup A_2$与$A_3\cup A_4$相互独立。
\item \hyperref[ex:w04-h4]{作业第4题：失败无穷多次}：$\P(A)=1$。
\end{enumerate}
\subsection*{第5周}
\begin{enumerate}
\item \hyperref[ex:w05-t33-4]{教材第33页第4题}：$2/5$。
\item \hyperref[ex:w05-t44-5]{教材第44页第5题}：$\P(a\leq\xi\leq b)\geq1-\alpha-\beta$。
\item \hyperref[ex:w05-h2]{作业第2题：区间投点的重复试验}：(1) $1-(2/3)^k$、$1-(1/2)^k$；(2) $1-(1/4)^k$；(3) 不独立，可改取$B'=[1/6,2/3]$；(4) $C,D$独立当且仅当$A,B$独立；(5) $1/3$。
\item \hyperref[ex:w05-h3]{作业第3题：示性函数的乘积与并}：$I_{AB}=I_AI_B$，$I_{A\cup B}=I_A+I_B-I_AI_B$，均为逐点恒等式。
\item \hyperref[ex:w05-h4]{作业第4题：随机变量与常数取最小值}：$Y=X\wedge a$是原概率空间上的随机变量。
\end{enumerate}
\subsection*{第6周}
\begin{enumerate}
\item \hyperref[ex:w06-h1]{作业第1题：取出号码的中间值}：$X$在2、3、4处的概率分别为$3/10,2/5,3/10$，分布函数如上。
\item \hyperref[ex:w06-h2]{作业第2题：随机选盒后取球}：(1) $1/6,11/18,2/9$；(2) 分布函数如上。
\item \hyperref[ex:w06-h3]{作业第3题：分段分布函数的跳跃与连续部分}：(1) $1/10,1/5,7/60,1/12$，四点总概率为$1/2$；(2) $1/10,2/5,0$。
\item \hyperref[ex:w06-h4]{作业第4题：由单点概率确定分布函数参数}：$A=1/2$，$B=1/10$；所求概率为$1/4$。
\item \hyperref[ex:w06-h5]{作业第5题（思考题，选做）：分布函数的一致收敛}：$\displaystyle\sup_{t\in\R}|F_n(t)-F(t)|\to0$。
\end{enumerate}
\subsection*{第7周}
\begin{enumerate}
\item \hyperref[ex:w07-b52-1]{教材第52页第1题}：已通过的路口数满足 $\P(X=k)=2^{-(k+1)}$，$k\geq0$；包含当前路口的计数服从参数 $1/2$ 的几何分布。
\item \hyperref[ex:w07-b52-2]{教材第52页第2题}：$0<p<1$ 时，$\P(\xi=k)=p^k(1-p)+(1-p)^kp$，$k\geq1$；$p=0,1$ 时游程无限长。
\item \hyperref[ex:w07-b53-3]{教材第53页第3题}：$\P(X=k)=(1-p)^{k-1}p$（$1\leq k\leq4$），$\P(X=5)=(1-p)^4$。
\item \hyperref[ex:w07-b53-4]{教材第53页第4题}：$\eta\sim P(\lambda p)$，即 $\P(\eta=k)=e^{-\lambda p}(\lambda p)^k/k!$。
\item \hyperref[ex:w07-b53-5]{教材第53页第5题}：奇数概率为 $(1-e^{-2\lambda})/2$，偶数概率为 $(1+e^{-2\lambda})/2$。
\item \hyperref[ex:w07-b58-3]{教材第58页第3题}：$a=1/2$，$\P(X<3/2)=7/8$。
\item \hyperref[ex:w07-h2]{作业第2题：三角形密度的区间概率}：$59/72$。
\item \hyperref[ex:w07-h3]{作业第3题：储油量}：至少 $61$ 千升。
\item \hyperref[ex:w07-h4]{作业第4题：由归一化求密度参数}：$C=1/4$，所求概率为 $1-\tfrac32e^{-1/2}$。
\end{enumerate}
\subsection*{第8周}
\begin{enumerate}
\item \hyperref[ex:w08-b57-2]{教材第57页第2题}：$pf+qg$ 是密度；其分布函数为 $pF_f+qF_g$，其中 $F_f,F_g$ 分别是 $f,g$ 对应的分布函数。
\item \hyperref[ex:w08-b58-4]{教材第58页第4题}：$F(t)=\exp(-e^{-t})$，$t\in\R$。
\item \hyperref[ex:w08-h2]{作业第2题：由三角形密度求分布函数}：\[
F(t)=\begin{cases}0,&t<0,\\t^2/2,&0\leq t\leq1,\\1-(2-t)^2/2,&1<t<2,\\1,&t\geq2.\end{cases}
\]
\item \hyperref[ex:w08-h3]{作业第3题：对称密度的分布函数}：$F(c-t)+F(c+t)=1$ 对全部 $t\in\R$ 成立；特别地 $F(c)=1/2$。
\item \hyperref[ex:w08-h4]{作业第4题：先到列车的方向}：两问均为 $2/3$。
\item \hyperref[ex:w08-h5]{作业第5题：修理时间}：第一问 $0.368$；第二问字面“恰好10小时”为 $0.000$，若指10小时内修好则为 $0.393$。
\item \hyperref[ex:w08-h6]{作业第6题：旧车的剩余里程}：指数分布下为 $0.368$；保留题面“10公里”并统一单位后，均匀分布下为 $0.500$；若原意为“10千里”，则为 $0.333$。
\item \hyperref[ex:w08-h7]{作业第7题：正态分布的中心}：$\mu=3$。
\end{enumerate}
\subsection*{第9周}
\begin{enumerate}
\item \hyperref[ex:w09-b58-4]{教材第58页第4题（再次布置）}：$F(t)=\exp(-e^{-t})$，$t\in\R$。
\item \hyperref[ex:w09-b58-5]{教材第58页第5题}：$X\sim N(a,\sigma^2)$，密度如上。
\item \hyperref[ex:w09-b62-1]{教材第62页第1题}：$1/3$。
\item \hyperref[ex:w09-h2]{作业第2题：正态分布的区间概率}：依次为 $0.7977$、$0.6827$、$0.8413$。
\item \hyperref[ex:w09-h3]{作业第3题：由分位点确定收入分布}：$\mu=250000$、$\sigma\approx103782.1553$；两问概率分别为 $0.3150$、$0.1363$。
\item \hyperref[ex:w09-h4]{作业第4题：正态分位点}：$c\approx14.56$。
\item \hyperref[ex:w09-h5]{作业第5题：线段分割的长短比}：$2/5$。
\item \hyperref[ex:w09-h6]{作业第6题：由均匀分布构造指数分布}：$Y\sim\operatorname{Exp}(\lambda)$，密度为 $\lambda e^{-\lambda y}$（$y>0$），其他处取 $0$。
\item \hyperref[ex:w09-h7]{作业第7题：指数变量的尺度变换}：$Y\sim\operatorname{Exp}(\lambda/c)$。
\item \hyperref[ex:w09-h8]{作业第8题：均匀变量的绝对值}：$Y\sim U(0,1/2)$。
\item \hyperref[ex:w09-h9]{作业第9题：指数变量的余弦}：\[
f_Y(y)=\begin{cases}
\dfrac{e^{-\arccos y}+e^{-2\pi+\arccos y}}
{(1-e^{-2\pi})\sqrt{1-y^2}},&-1<y<1,\\[6pt]
0,&\text{其他。}
\end{cases}
\]
\item \hyperref[ex:w09-h10]{作业第10题及思考题：概率积分变换}：$F(X)\sim U(0,1)$；即使 $F$ 有平坦区间，结论仍成立。这里保留了题设中 $X$ 连续型的条件。
\end{enumerate}
\subsection*{第10周}
\begin{enumerate}
\item \hyperref[ex:w10-b72-1]{教材第72页第1题}：联合分布为 $p(i,j)=1/(nj)$（$1\leq i\leq j\leq n$）；两个边缘分布如上。
\item \hyperref[ex:w10-b72-2]{教材第72页第2题}：$P(\xi=\eta)=0$。
\item \hyperref[ex:w10-b72-4]{教材第72页第4题}：$a=2$；$f_\xi(x)=2e^{-2x}I_{\{x>0\}}$；后两问为 $(1-e^{-4})(1-e^{-1})$、$(1-e^{-2})^2$。
\item \hyperref[ex:w10-b72-5]{教材第72页第5题}：$f_\xi(x)=3(1-x)^2I_{\{0<x<1\}}$，$f_\eta(y)=(3y-\frac32y^2)I_{\{0<y<1\}}$。
\item \hyperref[ex:w10-b72-6]{教材第72页第6题}：不是二维概率分布函数。
\item \hyperref[ex:w10-b81-4]{教材第81页第4题}：$f_\xi(x)=2xI_{\{0<x<1\}}$、$f_\eta(y)=2yI_{\{0<y<1\}}$，且二者独立。
\item \hyperref[ex:w10-h2]{作业第2题：由事件独立性确定联合分布}：$a=0.4,b=0.1$；$X$ 的分布为 $(0.8,0.2)$，$Y$ 的分布为 $(0.5,0.5)$；相互独立。
\item \hyperref[ex:w10-h3]{作业第3题：三项分布的边缘分布}：$X_1\sim B(n,p_1)$，$X_2\sim B(n,p_2)$，分布列如上。
\item \hyperref[ex:w10-h4]{作业第4题：双曲线之间的联合密度}：$A=3/2$；边缘密度如上；两概率为 $5/16$、$19/32$。
\end{enumerate}
\subsection*{第11周}
\begin{enumerate}
\item \hyperref[ex:w11-b80-1]{教材第80页第1题}：$f_{\xi\mid\eta}(x\mid y)=\sin(x+y)/(\sin y+\cos y)$，$0<x<\pi/2$；区间外为 $0$。
\item \hyperref[ex:w11-b80-2]{教材第80页第2题}：联合分布列如上，支持为 $n\geq1,0\leq k\leq n$ 的整数对。
\item \hyperref[ex:w11-b81-3]{教材第81页第3题}：联合密度如上。
\item \hyperref[ex:w11-b81-5]{教材第81页第5题}：$f_\xi(x)=4x(1-x^2)I_{\{0<x<1\}}$，$f_\eta(y)=4y^3I_{\{0<y<1\}}$；不独立。
\item \hyperref[ex:w11-h2]{作业第2题：中间次序与两个变量之和}：（1）$1/3$；（2）
\[
f_S(s)=\begin{cases}\frac23s^3,&0<s\leq1,\\-\frac23s^3+4s-\frac83,&1<s<2,\\0,&\text{其他}.\end{cases}
\]
\item \hyperref[ex:w11-h3]{作业第3题：首次非零结果与等待次数}：$N\sim G(1-p_0)$，$P(X=j)=p_j/(1-p_0)$，且 $N,X$ 独立。
\item \hyperref[ex:w11-h4]{作业第4题：保龄球得分}：（1）$0.6554$；（2）$0.2119$。
\item \hyperref[ex:w11-h5]{作业第5题：正态分布的等尾概率}：$x=20+5\sqrt2$。
\item \hyperref[ex:w11-h6]{作业第6题：逐次选数的条件分布}：联合分布为 $1/(5k)$（$1\leq j\leq k\leq5$）；条件分布如表；$X,Y$ 不独立。
\item \hyperref[ex:w11-h7]{作业第7题：三角形区域上的条件密度}：条件密度如上；条件概率为 $5/12$。
\item \hyperref[ex:w11-h8]{作业第8题：对称三角形区域上的条件密度}：条件密度如上；两个条件概率为 $1/6$、$20/27$。
\end{enumerate}
\subsection*{第12周}
\begin{enumerate}
\item \hyperref[ex:w12-b90-1]{教材第90页第1题}：$\xi+\eta$ 服从上述三角形密度。
\item \hyperref[ex:w12-b90-3]{教材第90页第3题}：$S\sim P(\lambda_1+\lambda_2)$，$\xi_1\mid(S=n)\sim B(n,\lambda_1/(\lambda_1+\lambda_2))$。
\item \hyperref[ex:w12-b90-4]{教材第90页第4题}：$\xi_1+\xi_2$ 与 $\xi_1/\xi_2$ 独立。
\item \hyperref[ex:w12-b90-6]{教材第90页第6题}：$\eta_n\sim\operatorname{Exp}(\lambda_1+\cdots+\lambda_n)$。
\item \hyperref[ex:w12-h2]{作业第2题：Gamma--泊松条件分布}：$W\mid(N=n)\sim\Gamma(t+n,\beta+1)$。
\item \hyperref[ex:w12-h3]{作业第3题：差、积、商的密度}：\[
f_{Y-X}(z)=\frac32(1-z^2)I_{\{0<z<1\}},\quad
f_{XY}(z)=3(1-\sqrt z)I_{\{0<z<1\}},\quad
f_{X/Y}(z)=I_{\{0<z<1\}}.
\]
\item \hyperref[ex:w12-h4]{作业第4题：指数样本的极差}：$n\geq2$ 时，$f_R(r)=(n-1)\lambda e^{-\lambda r}(1-e^{-\lambda r})^{n-2}I_{\{r>0\}}$；$n=1$ 时 $R=0$ 几乎必然。
\item \hyperref[ex:w12-h5]{作业第5题：五台发动机组成的机器寿命}：机器寿命为 $T_{(3)}$，密度为
\[
f_T(t)=30t(1+t)^2e^{-3t}[1-(1+t)e^{-t}]^2I_{\{t>0\}}.
\]
\item \hyperref[ex:w12-h6]{作业第6题：平方和与比值的联合分布}：$U\sim\operatorname{Exp}(1/2)=\chi^2(2)$，$V$ 服从标准柯西分布；联合密度如上，二者独立。
\item \hyperref[ex:w12-h7]{作业第7题：不同指数变量的最先发生概率}：第 $i$ 个变量最先发生的概率为 $\lambda_i/(\lambda_1+\cdots+\lambda_n)$。
\end{enumerate}
